\chapter{Results}
\label{ch:results}

This chapter reports what the simulator produced. It begins with the validity check that the
harness reproduces the expected sign, summarises all six sweeps in one table, and then takes
them in turn: the disclosure dial $B$, the responder count $n$, responder signal quality $q$,
signal dependence $\rho$, forecaster accuracy $q_F$ and the public prior $\pi$. It closes with
what the information measure can and cannot contribute.

All numbers come from a fixed seed at $200{,}000$ draws per cell, with the baseline
$n=5$, $q=0.65$, $\rho=0$, $q_F=0.75$, $B=8$, $\pi=0.5$ and one parameter varied at a time.
The baseline was fixed before any sweep was run.

\paragraph{How differences are judged.} Two intervals are used and they answer different
questions. A cell's own $95\%$ interval on $L$ is a paired difference between the two arms
on identical draws; at the baseline it is $\pm 0.0007$. That interval answers whether $L$
differs from zero in that cell. It is the wrong interval for asking whether $L$ differs
between two cells, because every cell in a sweep is run from the same seed and therefore
shares the same realised event and the same forecaster signal, so a between-cell difference
is itself paired. Its interval is typically an order of magnitude tighter, and the sweeps
record it for every step. Claims below about whether a change is resolved use the
between-cell interval; where a step is not resolved even by that test, it is reported as
unresolved.

\section{Validity check: the sign of the effect}
\label{sec:validity}

At the baseline the forecaster's expected surplus is $0.3749$ per unit when responders
cannot condition on the request and $0.1000$ when they can, so $L = 0.2749 \pm 0.0007$.

As Section~\ref{sec:cannot} set out, this is a check on the harness rather than a finding.
The standard adverse-selection result predicts the sign \cite{glosten1985bidask,
kyle1985continuous}, and a simulator producing $L \le 0$ would be evidence of an
implementation error. The check passes. The magnitude at this configuration is reported in
Section~\ref{sec:res-summary} rather than here, because it is conditional on parameters the
next sections vary.

One expectation should be lowered before the sweeps rather than after. Three of the six
sweeps produce behaviour that follows from the same adverse-selection logic that predicts
the sign; what they add is magnitude and shape rather than direction. Two sweeps, in $q$ and
in $\pi$, produce non-monotonicity that the sign alone does not anticipate, and the $q_F$
sweep produces the result with the clearest bearing on the report's argument.

\section{All six sweeps}
\label{sec:res-table}

\begin{table}[htbp]
\centering
\small
\begin{tabularx}{\textwidth}{llrrrrX}
\toprule
\textbf{Sweep} & \textbf{Range} & \multicolumn{2}{c}{\textbf{$L$}} & \multicolumn{2}{c}{\textbf{$I$ (nats)}} & \textbf{Shape} \\
 & & low & high & low & high & \\
\midrule
$B$    & 1 to 64        & 0.2615 & 0.2752 & 0.119 & 0.171 & rises, flat beyond $B=8$ \\
$n$    & 1 to 30        & 0.2273 & 0.2844 & \multicolumn{2}{c}{0.170 (invariant)} & rises, flat beyond $n{\approx}25$ \\
$q$    & 0.51 to 0.95   & 0.0662 & 0.2749 & 0.040 & 0.185 & peaks at $q=0.65$ \\
$\rho$ & 0 to 1         & 0.2274 & 0.2749 & \multicolumn{2}{c}{0.170 (invariant)} & falls throughout \\
$q_F$  & 0.55 to 0.95   & 0.0472 & 0.5441 & 0.007 & 0.520 & rises throughout \\
$\pi$  & 0.1 to 0.9     & 0.0411 & 0.2749 & 0.066 & 0.170 & peaks at $\pi=0.5$ \\
\bottomrule
\end{tabularx}
\caption{Endpoints of every sweep. The $L$ columns are the minimum and maximum over the
grid, not the values at the ends of the range, because two sweeps are non-monotonic. $I$ is
constant by construction in the $n$ and $\rho$ sweeps, for the reason given in
Section~\ref{sec:res-info}; it is a measured quantity only in the other four.}
\label{tab:results}
\end{table}

Table~\ref{tab:results} gives every endpoint, so the sections below can carry interpretation
rather than arithmetic. The grids are: $n \in \{1,2,3,4,5,6,8,10,12,15,20,25,30\}$;
$q$ and $q_F$ in steps of $0.05$ over their ranges; $\rho$ and $\pi$ in steps of $0.1$;
and $B \in \{1,2,3,4,6,8,12,16,24,32,48,64\}$.

\section{Disclosure granularity}
\label{sec:res-buckets}

\begin{figure}[htbp]
\centering
\includegraphics[width=\textwidth]{sweep_buckets.pdf}
\caption{Disclosure granularity $B$, on a logarithmic axis. Left: expected forecaster
surplus per unit in both arms; the no-disclosure arm is flat at $0.3749$ by construction,
because $B$ affects only what is disclosed. Right: the transfer $L$, solid line with diamonds and
$95\%$ intervals, on the left axis; the information $I$ in nats, dash-dot line with
triangles, on the right axis; both axes are anchored at zero. Almost the whole effect is present at $B=1$, where the request carries direction
alone.}
\label{fig:buckets}
\end{figure}

Figure~\ref{fig:buckets} sweeps the disclosure dial from one band to sixty-four. The
no-disclosure arm is flat at $0.3749$ across every value of $B$, which confirms that $B$
varies disclosure and nothing else.

The finding is how little the rest of the dial matters. At $B=1$, where the request reveals
only a direction, $L = 0.2615$. At $B=64$, where it reveals confidence to within a
sixty-fourth, $L = 0.2752$. Direction alone therefore accounts for $95.0\%$ of the total
transfer, and the entire range of confidence granularity accounts for the remaining $5\%$.

The curve flattens quickly. Every step up to $B=8$ is resolved; every step from $B=8$
onwards is not, on the between-cell test. The rise from $B=8$ to $B=64$ is $+0.00024$, which
is resolved when compared directly across the whole interval but is economically negligible
and is not attributable to any single step. The practical reading is that granularity beyond
about eight bands adds nothing worth having.

The information measure behaves differently from the surplus. $I$ rises from $0.119$ nats at
$B=1$ to $0.171$ at $B=64$, so finer disclosure teaches a responder appreciably more about
the event while moving very little additional surplus. The two quantities are in different
units and no ratio between them is meaningful, so the comparison here is directional only:
$I$ continues to rise over a range where $L$ has stopped.

\section{Responder count}
\label{sec:res-n}

\begin{figure}[htbp]
\centering
\includegraphics[width=\textwidth]{sweep_n.pdf}
\caption{Responder count $n$. Axes as in Figure~\ref{fig:buckets}. Competition raises the
forecaster's surplus in both arms but raises it less under disclosure, so the transfer grows
with $n$. The information series is flat because $I$ cannot depend on $n$; see
Section~\ref{sec:res-info}.}
\label{fig:n}
\end{figure}

Figure~\ref{fig:n} sweeps the number of responders from one to thirty. Whether more
responders help or hurt an informed forecaster is not settled by the sign of
the effect, because more responders means both more parties observing the request and more
competition for the business. The measured answer is that competition never dominates. $L$
rises from $0.2273$ at $n=1$ to $0.2844$ at $n=30$, and every step is resolved except the
last, from $n=25$ to $n=30$, which is $+7\times10^{-7}$ against an interval of
$\pm 1.3\times10^{-6}$.

The mechanism is visible in the two arms separately. Competition helps the forecaster in
both: surplus rises from $0.0016$ to $0.1171$ under disclosure and from $0.2289$ to $0.4015$
without it. It helps more in absolute terms in the no-disclosure arm, so the gap widens.
Adding responders to a disclosing mechanism improves the forecaster's position, and worsens
it relative to the same mechanism without disclosure. Both are true. This report treats the
relative comparison as the design-relevant one, because it is the only one that isolates the
disclosure decision from everything else the mechanism does, and that is an argument rather
than a measurement.

At $n=1$ the disclosure arm leaves the forecaster with $0.0016$ per unit, approximately zero.
This is not monopoly pricing: the single responder quotes its own posterior truthfully,
exactly as at every other $n$. What changes with $n$ is only that the forecaster
selects the most favourable of $n$ quotes, so at $n=1$ there is no order statistic to select
over, and a responder who has seen the direction holds a posterior close to the forecaster's
own.

\section{Responder signal quality}
\label{sec:res-q}

\begin{figure}[htbp]
\centering
\includegraphics[width=\textwidth]{sweep_q.pdf}
\caption{Responder signal quality $q$. Axes as in Figure~\ref{fig:buckets}. The transfer is
non-monotonic, peaking at $q = 0.65$ and collapsing as responders approach certainty, while
the information a responder gains falls throughout.}
\label{fig:q}
\end{figure}

Figure~\ref{fig:q} sweeps responder signal quality. $L$ is non-monotonic in $q$. It rises
from $0.2528$ at $q=0.51$ to a maximum of $0.2749$ at
$q=0.65$, then falls to $0.0662$ at $q=0.95$. The peak is resolved against both of its
neighbours: the step up from $q=0.60$ is $+0.0029$ and the step down to $q=0.70$ is
$-0.0048$, against between-cell intervals of $\pm 0.0001$ and $\pm 0.0002$. The peak is
modest in height, $8.7\%$ above the value at $q=0.51$, but the collapse on the right is
large: the transfer at $q=0.95$ is $24\%$ of its maximum.

The right-hand side is straightforward. A responder near $q=0.95$ already prices close to
the truth from its own signal, so learning the forecaster's direction shifts its quote very
little, and $I$ confirms this directly, falling from $0.185$ nats at $q=0.51$ to $0.040$ at
$q=0.95$.

The left-hand side needs more care, and the explanation is not the one that first suggests
itself. The forecaster's benchmark surplus is itself strongly non-monotonic in $q$, rising
from $0.2609$ at $q=0.51$ to a peak of $0.4123$ at $q=0.75$ before falling to $0.1777$ at
$q=0.95$. The quoted spread follows the same arc, and it is \emph{narrowest} at low $q$:
$0.019$ at $q=0.51$ against $0.404$ at $q=0.80$. Responders who know almost nothing all
quote close to the prior, so they quote close to each other, and the forecaster's
best-of-$n$ selection has almost nothing to select over. The benchmark surplus at low $q$ is
therefore small because responders are uniformly mispriced rather than dispersed, and $L$,
being a difference between two arms that are both small there, is bounded by how little
there was to take. $L$ peaks in the middle because that is where the benchmark surplus is
already substantial and the request still tells a responder something it does not know.

Note that $I$ is monotone in $q$ while $L$ is not. The information measure is consistent with
the right-hand collapse and cannot by itself account for the left-hand rise, which is
governed by the benchmark surplus and the dispersion of quotes rather than by what any one
responder learns.

\section{Dependence among responder signals}
\label{sec:res-rho}

\begin{figure}[htbp]
\centering
\includegraphics[width=\textwidth]{sweep_rho.pdf}
\caption{Signal dependence $\rho$. Axes as in Figure~\ref{fig:buckets}. The transfer falls
as responders become more dependent, but so does the forecaster's surplus in both arms.
Dependence is not good news for the forecaster: it shrinks the measured transfer by
shrinking what there was to take.}
\label{fig:rho}
\end{figure}

Figure~\ref{fig:rho} sweeps signal dependence from independent draws to a single shared
one. A note on the parameter first. $\rho$ is the probability that a responder copies a single common
draw rather than drawing its own, so it governs dependence among responder signals but is
not a correlation coefficient. Every signal is marginally correct with probability $q$ at
every $\rho$, and the signals are conditionally independent given the event at $\rho = 0$
rather than unconditionally independent; measured pairwise correlation at $\rho = 0$ is
$0.09$, not zero. Chapter~\ref{ch:method} describes the construction; the label should be
read as dependence rather than as a correlation.

$L$ falls monotonically as $\rho$ rises, from $0.2749$ at $\rho=0$ to $0.2274$ at $\rho=1$,
a reduction of $17\%$, and every step is resolved.

Reporting only that would mislead. The forecaster's surplus falls in \emph{both} arms as
$\rho$ rises: from $0.3749$ to $0.2288$ without disclosure, and from $0.1000$ to $0.0014$
with it. In proportional terms the disclosure arm is devastated, retaining $1.4\%$ of its
value at $\rho=0$. The quoted spread collapses to zero at $\rho=1$, because responders that
all copy the same draw quote identically and competition between them ceases to operate. At
$\rho=1$ the $n$ responders behave as one, which is why the numbers there closely match
those at $n=1$: $L = 0.2274$ against $0.2273$, and a disclosure-arm surplus of $0.0014$
against $0.0016$.

The mechanism is the destruction of competition, not a reduction in what any responder can
exploit; Section~\ref{sec:res-info} shows that what an individual responder learns from the
request does not change with $\rho$ at all. So dependence reduces the measured transfer
while making the forecaster worse off overall. $L$ is a difference, and a difference can
shrink because the mechanism became fairer or because there was less left to take. Here it
is the second, and any reading of the $\rho$ sweep that treats falling $L$ as an improvement
would be wrong. Chapter~\ref{ch:evaluation} takes this as a caution about the metric itself.

\section{Forecaster accuracy and the public prior}
\label{sec:res-qf}

Two parameters govern the setting rather than the responders. Both were swept because the
headline magnitude depends on them and neither is otherwise defensible as a fixed choice.

\begin{figure}[htbp]
\centering
\includegraphics[width=\textwidth]{sweep_qf.pdf}
\caption{Forecaster directional accuracy $q_F$. Axes as in Figure~\ref{fig:buckets}. Both
the transfer and the share of the forecaster's benchmark surplus that it represents rise
throughout: a better forecaster loses proportionally more to disclosure.}
\label{fig:qf}
\end{figure}

Figure~\ref{fig:qf} sweeps forecaster accuracy. $L$ rises monotonically and steeply in
$q_F$, from $0.0472$ at $q_F=0.55$ to $0.5441$ at
$q_F=0.95$, with every step resolved. The share of the forecaster's benchmark surplus that
does not survive disclosure rises with it, from $26.1\%$ at $q_F=0.55$ to $73.3\%$ at the
baseline $q_F=0.75$ and $95.9\%$ at $q_F=0.95$.

This is the sweep with the clearest bearing on the report's argument, and it is worth
stating carefully because it is easy to overstate. What the model shows is that under this
mechanism the proportion of an informed forecaster's advantage that transfers to responders
increases with how well informed the forecaster is. A forecaster who is barely better than
the prior keeps roughly three quarters of a small edge; a forecaster who is nearly always
right keeps roughly four per cent of a large one. Disclosure therefore bears hardest on
precisely the participants whose information the market exists to aggregate. Whether that
constitutes the transfer the report's definition is concerned with, or ordinary compensation
to responders for facing a better-informed counterparty, is the question
Chapter~\ref{ch:evaluation} takes up; the measurement does not settle it.

\begin{figure}[htbp]
\centering
\includegraphics[width=\textwidth]{sweep_pi.pdf}
\caption{Public prior $\pi$. Axes as in Figure~\ref{fig:buckets}. The transfer is symmetric
about an even prior and falls away sharply towards either certainty.}
\label{fig:pi}
\end{figure}

Figure~\ref{fig:pi} sweeps the public prior. $L$ is symmetric in $\pi$ and peaks at
$\pi = 0.5$, falling from $0.2749$ there to $0.0418$
at $\pi = 0.1$ and $0.0411$ at $\pi = 0.9$. Every step is resolved, and the near-equality of
the two tails is a useful internal consistency check, since nothing in the model privileges
either outcome. The reading is that disclosure matters most when the public estimate is
least informative. On a question the public prior already nearly settles, a forecaster's
direction is close to predictable and reveals little.

\section{What the information measure contributes}
\label{sec:res-info}

Reporting $I$ alongside $L$ was intended to link how much responders learn to how much
surplus moves. The link is real in some sweeps and absent by construction in others, and the
distinction matters.

$I$ is computed from a responder's own signal, the observed action, and the publicly known
parameters $q$, $\pi$, $q_F$ and $B$. It contains no $n$ and no $\rho$. Because every
responder signal is marginally correct with probability $q$ whatever $\rho$ is, and because
the action is independent of a responder's signal given the event, the joint distribution
that $I$ averages over does not depend on either parameter. \textbf{$I$ therefore cannot vary with
$n$ or with $\rho$.} The flatness observed in those two sweeps, at $0.170$ nats throughout,
is a check that the metric is implemented as defined; it is not a measured null and it is
not evidence about disclosure. It is reported as such in Table~\ref{tab:results}, and the
figures show a flat series against an axis anchored at zero so that the constancy is legible
rather than magnified into an apparent trend.

What follows from that invariance is nevertheless useful, and it is a logical consequence
rather than an empirical discovery: since $I$ is constant while $L$ varies by $25\%$ across
the $n$ sweep and $17\%$ across the $\rho$ sweep, $L$ is not a function of $I$ alone. What
converts a responder's learning into transferred surplus is the competitive structure, not
the learning by itself.

In the four sweeps where $I$ can move, it does. It falls monotonically in $q$, rises in $B$,
rises steeply in $q_F$, and peaks with $L$ at $\pi = 0.5$. In three of those four it moves
with $L$; in $q$ it does not, because $L$ is non-monotonic there and $I$ is not.

The consequence for Chapter~\ref{ch:evaluation} is narrow and should not be inflated. A
mechanism that limited how much a request reveals, measured in nats, would not thereby
determine the economic transfer, because the same $I$ is consistent with a range of $L$
spanning $25\%$. It remains true that $I = 0$ forces $L = 0$ in this model, so limiting
disclosure does constrain the transfer; the point is that it does not pin it down, and that
the competitive environment must be considered alongside any disclosure bound.

\section{Robustness to the seed}
\label{sec:res-seeds}

Every figure in this chapter is drawn from one realisation at the seed fixed in
Chapter~\ref{ch:method}. The paired intervals bound Monte Carlo error within that
realisation and say nothing about how far the results would move at a different one. The
sweeps were therefore re-run in full at five seeds. Table~\ref{tab:seeds} reports the result.

\begin{table}[htbp]
\centering
\small
\begin{tabularx}{\textwidth}{Xll}
\toprule
\textbf{Claim or statistic} & \textbf{Across five seeds} & \textbf{Spread} \\
\midrule
$L > 0$ at the baseline & holds at every seed & --- \\
$L$ increases monotonically in $n$ & holds at every seed & --- \\
$L$ decreases monotonically in $\rho$ & holds at every seed & --- \\
$L$ increases monotonically in $q_F$ & holds at every seed & --- \\
The $q$ sweep peaks at $q = 0.65$ & holds at every seed & --- \\
The $\pi$ sweep peaks at $\pi = 0.5$ & holds at every seed & --- \\
\midrule
Baseline $L$ & $0.2744$ to $0.2755$ & $0.0011$ \\
Share of benchmark surplus lost at baseline & $73.3\%$ to $73.8\%$ & $0.5$ pp \\
Direction's share of the transfer & $94.9\%$ to $95.2\%$ & $0.3$ pp \\
Share lost at $q_F = 0.55$ & $26.1\%$ to $26.6\%$ & $0.5$ pp \\
Share lost at $q_F = 0.95$ & $95.9\%$ to $96.0\%$ & $0.1$ pp \\
\bottomrule
\end{tabularx}
\caption{Seed robustness. Every qualitative claim in this chapter survives all five seeds,
including both locations of a non-monotonic peak, and no headline statistic moves by as much
as one percentage point.}
\label{tab:seeds}
\end{table}

Two observations are worth making. The first is that the two non-monotonic results, which
are the ones a single realisation could most plausibly have manufactured, put their peak at
the same grid point at every seed. The second is that the spread in baseline $L$ across
seeds, $0.0011$, is slightly larger than a single cell's $95\%$ interval of $\pm 0.0007$,
which is what one would expect of the range of five draws and is a mild reminder that the
within-seed interval is not a bound on total uncertainty.

This check bounds seed variation. It does not bound specification variation: every run uses
the same model, and the sensitivity of these results to the modelling choices set out in
Section~\ref{sec:cannot} is untested.

\section{What this chapter established}
\label{sec:res-summary}

Demonstrated, in the sense of measured by the experiment:

\begin{enumerate}
  \item Under a matched benchmark holding competition, signal quality, dependence and risk
    identical, letting responders condition on the request transfers
    $L = 0.2749 \pm 0.0007$ per unit at the baseline, $73.3\%$ of the forecaster's benchmark
    surplus. The sign was expected; the magnitude is a property of this model at this
    configuration, and Section~\ref{sec:res-qf} shows how strongly it depends on $q_F$.
  \item Direction alone accounts for $95.0\%$ of that transfer. Coarsening the disclosed
    confidence from sixty-four bands to one recovers only $5\%$ of it, and no step beyond
    $B=8$ is individually resolved.
  \item $L$ rises in the number of responders, with every step resolved except the last.
    Competition helps the forecaster in both arms but helps less under disclosure, so more
    competition does not compensate for disclosure. At $n=1$ the disclosure-arm surplus is
    approximately zero, through the absence of any quote to select over rather than through
    monopoly pricing.
  \item $L$ is non-monotonic in responder signal quality, peaking at $q=0.65$ and falling to
    $24\%$ of that maximum by $q=0.95$. The peak is resolved against both neighbours.
  \item $L$ falls as responder signals become more dependent, but the forecaster's surplus
    falls in both arms and the disclosure arm retains $1.4\%$ of its value at $\rho=1$. The
    metric improves while the forecaster's position worsens.
  \item The share of benchmark surplus lost to disclosure rises monotonically with
    forecaster accuracy, from $26.1\%$ at $q_F=0.55$ to $95.9\%$ at $q_F=0.95$.
  \item $L$ is symmetric in the public prior and peaks where the prior is least informative.
\end{enumerate}

Established by construction rather than measured: $I$ does not vary with $n$ or $\rho$, from
which it follows that $L$ is not a function of $I$ alone.

Not established by this chapter: any magnitude for a live venue; anything about accumulation
across repeated requests, which is not modelled; anything about multi-leg disclosure, which
was not implemented; and, most consequentially for what follows, any separation between surplus
transferred because of information and surplus that would compensate a responder for risk.
The model's responders are risk-neutral and quote their own posteriors without shading, so
that separation is imposed by construction rather than measured, and $L$ should not be read as
pure information rent. Chapter~\ref{ch:evaluation} asks how much of $L$ meets the definition
this report adopted, and how much of it is ordinary adverse selection.
